% Solutions to the problems of Chapter 6 (Sensor Fusion in Multi-Agent Systems). October 2026 revision.
\begin{sol}{prob:ch06:1}
For the running network, \eqref{eq:consensus-rate} gives $\rho=\max\{|1-0.3\cdot1.382|,|1-0.3\cdot4.618|\}=\max\{0.585,0.385\}=0.585$, and a reduction by a factor of $100$ requires $k\ge\ln 100/\ln(1/0.585)\approx8.6$, so $9$ rounds. The optimal step is $\epsilon^{\star}=2/(\lambda_2+\lambda_5)=2/6=1/3$, with $\rho^{\star}=(4.618-1.382)/6=0.539$ and $8$ rounds. For the path the eigenvalues are $0$, $0.382$, $1.382$, $2.618$, and $3.618$; with $\epsilon=0.3$, $\rho=\max\{0.885,0.085\}=0.885$ and $38$ rounds are needed, while $\epsilon^{\star}=2/(0.382+3.618)=0.5$ gives $\rho^{\star}=0.809$ and $22$ rounds. In both cases $\epsilon^{\star}=1/d_{\max}$ ($d_{\max}=3$ for the network, $2$ for the path) lies on the boundary of the sufficient condition $\epsilon<1/d_{\max}$, but the exact condition is $\epsilon<2/\lambda_N$, that is, $\epsilon<0.433$ and $\epsilon<0.553$, which is satisfied: the bound $\lambda_N\le2d_{\max}$ is not tight for these graphs. Removing two edges multiplies the number of rounds by about four, because $\lambda_2$ falls from $1.382$ to $0.382$.
\end{sol}

\begin{sol}{prob:ch06:2}
For an undirected graph $\mathbf{L}$ is symmetric and $\mathbf{L}\mathbf{1}=\mathbf{0}$, hence $\mathbf{1}^{\mathsf T}\mathbf{L}=\mathbf{0}^{\mathsf T}$ and $\mathbf{1}^{\mathsf T}\mathbf{z}(k+1)=\mathbf{1}^{\mathsf T}\mathbf{z}(k)$, which is \eqref{eq:consensus-average-preserved}. For the weighted iteration, $\mathbf{1}^{\mathsf T}\mathbf{W}\mathbf{z}(k)=\mathbf{1}^{\mathsf T}\mathbf{z}(k)$ for every $\mathbf{z}(k)$ if and only if $\mathbf{1}^{\mathsf T}\mathbf{W}=\mathbf{1}^{\mathsf T}$, i.e., if and only if the columns of $\mathbf{W}$ sum to one. The given matrix has column sums $5/6$, $4/3$, and $5/6$, so it is not column-stochastic. Solving $\boldsymbol{\pi}^{\mathsf T}\mathbf{W}=\boldsymbol{\pi}^{\mathsf T}$ gives $\pi_1=\pi_1/2+\pi_2/3$ and $\pi_3=\pi_3/2+\pi_2/3$, i.e., $\pi_1=\pi_3=2\pi_2/3$; normalising, $\boldsymbol{\pi}=[2/7\ \ 3/7\ \ 2/7]^{\mathsf T}$. Since the graph (an undirected path with asymmetric weights) is connected and $\mathbf{W}$ has positive diagonal, $\mathbf{W}^{k}\to\mathbf{1}\boldsymbol{\pi}^{\mathsf T}$, and the agents agree on $\boldsymbol{\pi}^{\mathsf T}\mathbf{z}(0)=(20+36+22)/7=78/7\approx11.14$ instead of the average $11$: the central agent, which listens to everybody, is over-represented.
\end{sol}

\begin{sol}{prob:ch06:3}
Scalar case: $P_{\mathrm{CI}}(\omega)^{-1}=\omega+(1-\omega)/4=(1+3\omega)/4$, so $P_{\mathrm{CI}}(\omega)=4/(1+3\omega)$ is decreasing in $\omega$ and minimal at $\omega^{\star}=1$ with $P_{\mathrm{CI}}=1$: CI simply keeps the better estimate and discards the worse one, because with unknown correlation the second estimate might be a degraded copy of the first and could add nothing. Two-dimensional case: $\mathbf{P}_{\mathrm{CI}}^{-1}=\diag\!\big(\omega+(1-\omega)/4,\ \omega/4+(1-\omega)\big)$, so $\tr\mathbf{P}_{\mathrm{CI}}=4/(1+3\omega)+4/(4-3\omega)$. Setting the derivative to zero, $(1+3\omega)^{2}=(4-3\omega)^{2}$, gives $\omega^{\star}=1/2$ and $\mathbf{P}_{\mathrm{CI}}=\diag(1.6,1.6)$. Independent fusion would claim $\mathbf{P}_{\mathrm{f}}=(\mathbf{P}_1^{-1}+\mathbf{P}_2^{-1})^{-1}=\diag(0.8,0.8)$, i.e., twice as much information along each axis; CI is conservative but direction-aware, since along each axis it never claims more information than the better of the two sources ($P=1$) and here, because of the symmetric weight, it claims less.
\end{sol}

\begin{sol}{prob:ch06:4}
With prior variance $1$ (information $1$) and one measurement of variance $1$ (information $1$), each node has posterior information $2$ and variance $1/2$. The naive sum gives information $4$, i.e., a claimed variance of $1/4$. The centralised filter counts the prior once and both measurements, information $3$, variance $1/3$. The naive estimate is $\widehat{x}=(\widehat{x}_1+\widehat{x}_2)/2$ with $\widehat{x}_i=(m+y_i)/2$, where $m$ is the prior mean; hence $\widehat{x}=m/2+(y_1+y_2)/4$ and its error has variance $1/4+2/16=3/8=0.375$. The expected NEES is the ratio of the actual to the claimed variance, $0.375/0.25=1.5$ instead of $1$, and the naive estimate is also less accurate than the centralised one ($0.375>0.333$): double counting the prior both overstates the confidence and misweights the data. In Example~\ref{ex:mas-distributed-tracking} the same mechanism, applied recursively with feedback, inflates the information at every step and the NEES grows without bound.
\end{sol}

\begin{sol}{prob:ch06:5}
On the ring-with-chord graph the Metropolis matrix has second largest eigenvalue modulus $0.6545$, so three rounds already reduce the disagreement by a factor of about $3.6$; the simulation gives a position RMSE of $0.96$\,m, $0.92$\,m, and $0.91$\,m, with mean NEES $2.05$, $2.02$, and $2.01$, for $L=1$, $3$, and $10$. On the path graph the second largest eigenvalue modulus is larger (the graph is less connected, $\lambda_2=2-2\cos(\pi/5)\approx0.38$ for the Laplacian), (second largest eigenvalue modulus $0.8727$ of the Metropolis matrix), so the improvement from $L=1$ to $L=10$ is larger: with $20$ Monte Carlo runs the RMSE decreases from $1.08$\,m to $0.97$\,m on the path against $1.01$\,m to $0.96$\,m on the ring with chord, while the mean NEES stays near $2$ in all cases. The script \codeentry{P6.5} produces the full table.
\end{sol}

\begin{sol}{prob:ch06:6}
Renormalising the Metropolis weights over the received messages keeps every round a convex combination of posteriors, which is still a covariance-intersection step; the CI covariance bound is preserved under the stated linear model and independent packet-loss assumptions if the inputs satisfy it. A mean NEES below $4$ is only a scalar diagnostic, and the realised RMSE need not increase monotonically with $p$ because fewer rounds are effective and the end nodes receive less information. Average preservation, however, is lost because the time-varying weight matrices are no longer doubly stochastic, so the nodes agree on slightly different values from run to run. With $20$ Monte Carlo runs and $L=3$ the position RMSE is $0.967$\,m, $0.964$\,m, and $1.043$\,m for $p=0$, $0.3$, and $0.6$, with mean NEES $1.92$, $1.73$, and $1.85$. The implementation keeps the original self-weight and divides all surviving weights by their row sum. The script \codeentry{P6.6} reports the numbers.
\end{sol}

\begin{sol}{prob:ch06:7}
Every connected graph is $1$-robust. For $S_1=\{1,2,3\}$ and $S_2=\{4,5\}$, node~$1$ has one neighbour outside $S_1$ (node~$5$), node~$2$ none, node~$3$ one (node~$4$), node~$4$ one outside $S_2$ (node~$3$), and node~$5$ one (node~$1$): no node has two neighbours outside its own subset, so the network is not $2$-robust, let alone $3$-robust as the stated sufficient W-MSR theorem for $F=1$ assumes (the singletons $\{4\}$ and $\{5\}$ already violate $3$-robustness, since both nodes have only two neighbours). In the simulation, node~$4$, whose neighbours are nodes~$3$ and $5$, discards both values at every round, the one being smaller and the other larger than its own, and stays at $15$\,m, while nodes~$1$, $2$, and $3$ agree on $11.5$\,m after about $10$ rounds: the compliant nodes never agree (their spread remains $3.5$\,m), which is the failure mode of W-MSR on a graph that is not sufficiently robust. The script \codeentry{P6.7} reproduces it.
\end{sol}
